\documentclass[11pt]{article}

\usepackage{amsmath, amssymb, amsfonts}
\usepackage{geometry}
\geometry{margin=1in}

\title{Supplementary Information: Information Distortion, Effective Sample Size, and Corrected Covariance in Sequential Biophysical Measurements}
\author{}
\date{}

\begin{document}

\maketitle

\section*{1. Log-Likelihood and Score Definitions}

Consider a trajectory of $T$ sequential measurement intervals (e.g., time steps in a patch-clamp recording or frames in a fluorescence trace). The total log-likelihood is defined as the sum of per-interval log-likelihoods:

\begin{equation}
\ell(\theta) = \log L(\theta) = \sum_{t=1}^T \ell_t(\theta).
\end{equation}

The total score $S(\theta)$, representing the gradient of the log-likelihood, is the sum of the per-interval scores $s_t(\theta)$:

\begin{equation}
S(\theta) = \nabla_\theta \ell(\theta) = \sum_{t=1}^T s_t(\theta), \qquad s_t(\theta) = \nabla_\theta \ell_t(\theta).
\end{equation}

Expectations and covariances are calculated over independent realizations of full trajectories generated under parameter $\theta$.

\section*{2. Model Sensitivity: Gaussian Fisher Information (GFI)}

At each interval, the predictive likelihood is approximated as Gaussian: $y_t \sim \mathcal{N}(\mu_t(\theta), \sigma_t^2(\theta))$. The total Gaussian Fisher Information $\mathbf{H}(\theta)$ represents the model's predicted curvature or sensitivity:

\begin{equation}
\mathbf{H}(\theta) = \sum_{t=1}^T \mathbb{E}\left[\mathbf{I}_t^{\mathrm{G}}(\theta)\right] = - \mathbb{E}\left[\nabla_\theta^2 \ell(\theta)\right].
\end{equation}

For a scalar variance $\sigma_t^2$, the per-interval information is calculated from the gradients of the mean and variance:

\begin{equation}
\mathbf I_t^{\mathrm{G}}(\theta) = \frac{1}{\sigma_t^2(\theta)} \nabla_\theta \mu_t(\theta) \nabla_\theta \mu_t(\theta)^\top + \frac{1}{2\sigma_t^4(\theta)} \nabla_\theta \sigma_t^2(\theta) \nabla_\theta \sigma_t^2(\theta)^\top.
\end{equation}

\section*{3. Score Fluctuation Matrices}

To diagnose model failure, we compare predicted curvature $\mathbf{H}$ against observed score fluctuations. We distinguish between local noise and trajectory-wide dependencies:

\begin{itemize}
    \item \textbf{Sample Score Covariance ($\mathbf{J}_{\mathrm{sample}}$):} This represents the sum of variances of individual intervals. It captures the ``local'' volatility of the data, as if measurements were independent:
    \begin{equation}
    \mathbf{J}_{\mathrm{sample}} = \sum_{t=1}^T \mathrm{Cov}(s_t).
    \end{equation}
    \item \textbf{Total Score Covariance ($\mathbf{J}_{\mathrm{total}}$):} The covariance of the total score sum. This captures the ``global'' volatility, including both local variance and all cross-interval temporal covariances:
    \begin{equation}
    \mathbf{J}_{\mathrm{total}} = \mathrm{Cov}(S(\theta)) = \mathbf{J}_{\mathrm{sample}} + 2\sum_{i<j}\mathrm{Cov}(s_i,s_j).
    \end{equation}
\end{itemize}

\section*{4. Derivation of the Distortion-Corrected Covariance (DCC)}

To find the corrected uncertainty for parameter estimates $\hat{\theta}$, we perform a first-order Taylor expansion of the score function around the true parameters $\theta_0$:
\begin{equation}
S(\hat{\theta}) \approx S(\theta_0) + \nabla_\theta S(\theta_0) (\hat{\theta} - \theta_0).
\end{equation}
At the maximum likelihood estimate, the score is zero ($S(\hat{\theta}) = 0$). Replacing the empirical Hessian with its expectation ($-\mathbf{H}$), we solve for the estimation error:
\begin{equation}
(\hat{\theta} - \theta_0) \approx \mathbf{H}^{-1} S(\theta_0).
\end{equation}
The true covariance of the parameter estimates, $\Sigma_{\mathrm{DCC}}$, is the expectation of the squared error:
\begin{equation}
\Sigma_{\mathrm{DCC}} = \mathbb{E}\left[(\hat{\theta} - \theta_0)(\hat{\theta} - \theta_0)^\top\right] \approx \mathbf{H}^{-1} \mathbb{E}[S S^\top] \mathbf{H}^{-1}.
\end{equation}
Given $\mathbf{J}_{\mathrm{total}} = \mathrm{Cov}(S(\theta))$, we arrive at the corrected covariance:
\begin{equation}
\Sigma_{\mathrm{DCC}} = \mathbf{H}^{-1} \mathbf{J}_{\mathrm{total}} \mathbf{H}^{-1}.
\end{equation}
This correction ensures that if the actual score fluctuations are larger than the model-predicted curvature, the parameter uncertainty is appropriately inflated.

\section*{5. Information Distortion Analysis}

The total Information Distortion Matrix ($\mathbf{C}$) quantifies the discrepancy between actual fluctuations and the model-predicted curvature:
\begin{equation}
\mathbf{C}(\theta) = \mathbf{H}^{-1/2} \mathbf{J}_{\mathrm{total}} \mathbf{H}^{-1/2}.
\end{equation}
This distortion is decomposed into two physically distinct mechanisms:

\subsection*{5.1 Sample Distortion ($\mathbf{C}_{\mathrm{sample}}$)}
This captures the mismatch between per-interval score variance and Gaussian Fisher curvature, representing purely geometric modeling errors (e.g., non-Gaussianity):
\begin{equation}
\mathbf{C}_{\mathrm{sample}} = \mathbf{H}^{-1/2} \mathbf{J}_{\mathrm{sample}} \mathbf{H}^{-1/2}.
\end{equation}

\subsection*{5.2 Correlation Distortion ($\mathbf{R}$)}
This measures the inflation of uncertainty due to temporal dependencies (e.g., memory effects or slow kinetics):
\begin{equation}
\mathbf{R} = \mathbf{J}_{\mathrm{sample}}^{-1/2} \mathbf{J}_{\mathrm{total}} \mathbf{J}_{\mathrm{sample}}^{-1/2}.
\end{equation}
If measurements are temporally uncorrelated, $\mathbf{R} = \mathbf{I}$. The total distortion is the product of these two effects:
\begin{equation}
\mathbf{C} = \mathbf{C}_{\mathrm{sample}}^{1/2} \mathbf{R} \mathbf{C}_{\mathrm{sample}}^{1/2}.
\end{equation}

\section*{6. Effective Sample Size and Bias}

Together, these metrics distinguish between different failure modes of a biophysical model:

\begin{itemize}
    \item \textbf{Effective Sample Size ($T_{\mathrm{eff}}$):} If distortion is purely due to isotropic temporal correlation ($\mathbf{C} \approx \kappa \mathbf{I}$), the effective sample size is $T_{\mathrm{eff}} = T/\kappa$.
    \item \textbf{Distortion-Induced Bias (DIB):} If the expected score $\mathbf{g} = \mathbb{E}[S(\theta)] \neq 0$, it indicates a systematic drift or bias in the model:
    \begin{equation}
    \mathbf{b}_{\mathrm{DIB}} = \mathbf{H}^{-1} \mathbf{g}.
    \end{equation}
\end{itemize}

\end{document}